\documentclass{book}
%\documentclass{book}
%If you're going to print this, perhaps get ride of oneside so that \parts{} will have their own page!
\title{
	\begin{huge}
		\bf{automorphism representable by reals}
	\end{huge}
	}
\author{Nathanael Chwojko-Srawley}

%\date{} % If you want to specify a date

%  ╔══════════════════════════════════════════════════════════╗
%  ║                         packages                         ║
%  ╚══════════════════════════════════════════════════════════╝

\usepackage[a4paper, total={6in, 8in}]{geometry}
\usepackage[answerdelayed]{exercise}
\usepackage[font=itshape]{quoting}
\usepackage[most]{tcolorbox}         % Makes nice boxes for thms, cor, prop, or anything else.
\usepackage{amsmath, amssymb, amsthm} % Used to write better mathematical expressions (ex. \mathbb{R})
\usepackage{array}                   % \begin{table}{>{\em}c c} -> 1st column italic (bfseries for bold)
\usepackage{bbm}                     % for mathbbm{1}
\usepackage{blindtext}
\usepackage{commutative-diagrams}
\usepackage{dsfont}
\usepackage{dynkin-diagrams}         % for Dynkin diagrams
\usepackage{enumitem}
\usepackage{etoolbox}
\usepackage{euscript}
\usepackage{extarrows}               % Used to compile any UTF-8 Character (and supports Unicode)
\usepackage{fancyhdr}
\usepackage{float}                   % package with ``H'' for figures and tables, ex. Images, tables, etc
\usepackage{graphicx}                % Let's you use images in LaTeX; ex the \includegraphics command.
\usepackage{hhline}
\usepackage{hyperref}
\usepackage{cleveref}
\usepackage{inputenc}                % Used to compile any UTF-8 Character (and supports Unicode)
\usepackage{listings}                % Create nice colour code blocks (customized in preamble)
\usepackage{longtable}               % create tables that extend past one page
\usepackage{makeidx}
\usepackage{mathrsfs}
\usepackage{mathtools}
% \usepackage{microtype}               % makes nice text. Fixes some under-/over- filled box problems.
\usepackage{multicol}                % For multiple columns
\usepackage{pgfplots}
\usepackage{scalerel}
\usepackage{setspace}                % More flexibility on spacing in document.
\usepackage{stackengine}
\usepackage{stackrel}                % for left-right arrows, staking symbol above and bellow
\usepackage{thmtools}                % Customize theorem environment (in preamble).
\usepackage{tikz-cd}                 % for commutative diagrams https://tex.stackexchange.com/questions/419221/how-to-do-the-pushout-with-universal-property
\usepackage{tikz}                    % for commutative diagrams https://tex.stackexchange.com/questions/419221/how-to-do-the-pushout-with-universal-property
\usepackage{titlesec}
\usepackage{wasysym}                 % emoji's!
\usepackage{wrapfig}                 % To wrap text around figure
\usepackage{xcolor}                  % Colour text.

% \usepackage{comment}

% \usepackage{CJKutf8} % for some chinese

% \usepackage{dutchcal} % for lower case mathcal
% \pdfsuppresswarningpagegroup=1

% ╔═════════════════════════════════════════════════════╗
% ║           for figures via inkspace                  ║
% ╚═════════════════════════════════════════════════════╝

% to import figures via: https://github.com/gillescastel/inkscape-figures
\usepackage{import}
\usepackage{pdfpages}
\usepackage{transparent}


% This command is the ONLY command imported via the packages snippet, think of it as a tiny package written out
\newcommand{\incfig}[2][1]{%
    \def\svgwidth{#1\columnwidth}
    \import{./figures/}{#2.pdf_tex}
}

% ╔═══════════════════════════════════════════════════════════╗
% ║                 bibliography and citing                   ║
% ╚═══════════════════════════════════════════════════════════╝

\usepackage{csquotes}
% \usepackage[backend=biber,citestyle=alphabetic,bibstyle=authortitle]{biblatex}
\usepackage[backend=biber]{biblatex}

\addbibresource{bibliography.bib}

% ╔═══════════════════════════════════════════════════════════╗
% ║                    colours (colors)                       ║
% ╚═══════════════════════════════════════════════════════════╝

\definecolor{dkgreen}{rgb}{0,0.6,0}
\definecolor{gray}{rgb}{0.5,0.5,0.5}
\definecolor{mauve}{rgb}{0.58,0,0.82}
\definecolor{lightyellow}{rgb}{1,1,0.6}
\definecolor{reddish}{HTML}{A01A3A}

%  ╔══════════════════════════════════════════════════════════╗
%  ║               title/heading/footer format                ║
%  ╚══════════════════════════════════════════════════════════╝

\pagestyle{fancy}
\setlength\headheight{20pt}
\renewcommand{\sectionmark}[1]{\markright{\thesection.\ #1}}
\renewcommand{\chaptermark}[1]{\markboth{#1}{}}
\fancyfoot[C,CO]{\textcolor{gray}{Nathanael Chwojko-Srawley}}
\fancyfoot[R,RO]{\thepage}{}

\titleformat
{\chapter} % command
[display] % shape
{\bfseries\Huge\itshape} % format
{\thechapter} % label
{0.5ex} % sep
{
    \rule{\textwidth}{1pt}
    \vspace{1ex}
    \centering
}
[
\vspace{-0.5ex}
\rule{\textwidth}{0.3pt}
]

\setlength{\parindent}{0pt} % don't like indentation infront of paragraphs
\setlength{\parskip}{1ex}   %so that there is still space between paragarphs


%  ╔══════════════════════════════════════════════════════════╗
%  ║                       Shortcuts                          ║
%  ╚══════════════════════════════════════════════════════════╝

%\ExerciseLevelInToc

% I like original mathcal better, but need lowercase.
\DeclareFontFamily{OT1}{pzc}{}
\DeclareFontShape{OT1}{pzc}{m}{it}{<-> s * [1.10] pzcmi7t}{}
\DeclareMathAlphabet{\mathpzc}{OT1}{pzc}{m}{it}


% mathbb
\newcommand{\A}{{\mathbb{A}}}
\newcommand{\C}{{\mathbb{C}}}
\newcommand{\F}{{\mathbb{F}}}
\newcommand{\N}{{\mathbb{N}}}
\newcommand{\Q}{{\mathbb{Q}}}
\newcommand{\R}{{\mathbb{R}}}
\newcommand{\V}{{\mathbb{V}}}
\newcommand{\Z}{{\mathbb{Z}}}
\newcommand{\BI}{{\mathbb{I}}}
\renewcommand{\H}{{\mathbb{H}}}
\renewcommand{\P}{{\mathbb{P}}}


% mathcal
% \newcommand{\co}{{\mathfrak{o}}} %mathcal{o} looks like wreath product, so I changed it to mathfrak
\newcommand{\co}{{\mathpzc{o}}}
\newcommand{\CA}{{\mathcal{A}}}
\newcommand{\CB}{{\mathcal{B}}}
\newcommand{\CC}{{\mathcal{C}}}
\newcommand{\CD}{{\mathcal{D}}}
\newcommand{\CF}{{\mathcal{F}}}
\newcommand{\CG}{{\mathcal{G}}}
\newcommand{\CH}{{\mathcal{H}}}
\newcommand{\CI}{{\mathcal{I}}}
\newcommand{\CK}{{\mathcal{K}}}
\newcommand{\CL}{{\mathcal{L}}}
\newcommand{\CM}{{\mathcal{M}}}
\newcommand{\CN}{{\mathcal{N}}}
\newcommand{\CO}{{\mathcal{O}}}
\newcommand{\CQ}{{\mathcal{Q}}}
\newcommand{\CR}{{\mathcal{R}}}
\newcommand{\CS}{{\mathcal{S}}}
\newcommand{\CT}{{\mathcal{T}}}
\newcommand{\CV}{{\mathcal{V}}}
\newcommand{\CZ}{{\mathcal{Z}}}
% EXCPETION:
\newcommand{\CP}{{\mathcal{P}}}
% \newcommand{\CP}{\mathbb{C}\mathrm{P}}
\newcommand{\RP}{\mathbb{R}\mathrm{P}}


%Euscript
\newcommand{\EB}{\EuScript{B}}
\newcommand{\EC}{{\EuScript{C}}}
\newcommand{\EF}{{\EuScript{F}}}
\newcommand{\EL}{{\EuScript{L}}}
\newcommand{\EO}{{\EuScript{O}}}


%mathfrak (lowercase)
\newcommand{\fa}{{\mathfrak{a}}}
\newcommand{\fb}{{\mathfrak{b}}}
\newcommand{\fc}{{\mathfrak{c}}}
\newcommand{\fd}{{\mathfrak{d}}}
\newcommand{\fm}{{\mathfrak{m}}}
\newcommand{\fn}{{\mathfrak{n}}}
\newcommand{\fo}{{\mathfrak{o}}}
\newcommand{\fp}{{\mathfrak{p}}}
\newcommand{\fq}{{\mathfrak{q}}}


%mathfrak (upper case)
\newcommand{\FA}{{\mathfrak{A}}}
\newcommand{\FB}{{\mathfrak{B}}}
\newcommand{\FC}{{\mathfrak{C}}}
\newcommand{\FD}{{\mathfrak{D}}}
\newcommand{\FK}{{\mathfrak{K}}}
\newcommand{\FM}{{\mathfrak{M}}}
\newcommand{\FN}{{\mathfrak{N}}}
\newcommand{\FO}{{\mathfrak{O}}}
\newcommand{\FP}{{\mathfrak{P}}}


%mathscr
\newcommand{\ff}{{\mathscr{F}}}
\newcommand{\SA}{\mathscr{A}}
\newcommand{\SC}{\mathscr{C}}
\newcommand{\SF}{\mathscr{F}}
\newcommand{\SG}{\mathscr{G}}
\newcommand{\SH}{\mathscr{H}}
\newcommand{\SI}{\mathscr{I}}
\newcommand{\SK}{\mathscr{K}}
\newcommand{\SN}{\mathscr{N}}
\newcommand{\SQ}{\mathscr{Q}}
\newcommand{\SR}{\mathscr{R}}
\newcommand{\SV}{\mathscr{V}}
\newcommand{\SZ}{\mathscr{Z}}


% operatorname -- 2 letters (lowercase and upper case)
\newcommand{\ab}{{\operatorname{ab}}}
\newcommand{\ad}{{\operatorname{ad}}}
\newcommand{\ch}{{\operatorname{ch}}}
\newcommand{\ev}{{\operatorname{ev}}}
\newcommand{\id}{{\operatorname{id}}}
\newcommand{\im}{{\operatorname{im}}}
\newcommand{\nm}{{\operatorname{Nm}}}
\newcommand{\op}{{\operatorname{op}}}
\newcommand{\rk}{{\operatorname{rk}}}
\newcommand{\sh}{{\operatorname{sh}}}
\newcommand{\tr}{{\operatorname{tr}}}

\newcommand{\Ab}{{\operatorname{Ab}}}
\newcommand{\Br}{{\operatorname{Br}}}
\newcommand{\Ch}{{\operatorname{Ch}}}
\newcommand{\Cl}{{\operatorname{Cl}}}
\newcommand{\GL}{{\operatorname{GL}}}
\newcommand{\Nm}{{\operatorname{Nm}}}
\newcommand{\SL}{{\operatorname{SL}}}

\renewcommand{\Re}{{\operatorname{Re}}}
\renewcommand{\Im}{{\operatorname{Im}}}


%categories
\newcommand{\Grp}{{\textbf{Grp}}}
\newcommand{\Fld}{{\textbf{Fld}}}
\newcommand{\Sh}{{\textbf{Sh}}}
\newcommand{\Set}{{\textbf{Set}}}
\newcommand{\Mod}{{\textbf{Mod}}}
\newcommand{\PSh}{{\textbf{PSh}}}
\newcommand{\Sch}{{\textbf{Sch}}}
\newcommand{\Top}{{\textbf{Top}}}
\newcommand{\RgSp}{{\textbf{RgSp}}}
\newcommand{\Ring}{{\textbf{Ring}}}


%operatorname -- 3+ letters (lowercase and upper case)
\newcommand{\rad}{{\operatorname{rad}}}
\newcommand{\sep}{{\operatorname{sep}}}
\newcommand{\res}{{\operatorname{res}}}
\newcommand{\sgn}{{\operatorname{sgn}}}
\newcommand{\pre}{{\operatorname{pre}}}
\newcommand{\ord}{{\operatorname{ord}}}
\newcommand{\vol}{{\operatorname{vol}}}
\newcommand{\lcm}{{\operatorname{lcm}}}

\newcommand{\conj}{{\operatorname{conj}}}
\newcommand{\disc}{{\operatorname{disc}}}
\newcommand{\stab}{{\operatorname{stab}}}
\newcommand{\kdim}{{\operatorname{kdim}}}
\newcommand{\diag}{{\operatorname{diag}}}

\newcommand{\trdeg}{\operatorname{trdeg}}
\newcommand{\coker}{{\operatorname{coker}}}
\newcommand{\codim}{{\operatorname{codim}}}

\newcommand{\Ann}{{\operatorname{Ann}}}
\newcommand{\Aut}{{\operatorname{Aut}}}
\newcommand{\Der}{{\operatorname{Der}}}
\newcommand{\Div}{{\operatorname{Div}}}
\newcommand{\Emb}{{\operatorname{Emb}}}
\newcommand{\End}{{\operatorname{End}}}
\newcommand{\Ext}{{\operatorname{Ext}}}
\newcommand{\Gal}{{\operatorname{Gal}}}
\newcommand{\Hom}{{\operatorname{Hom}}}
\newcommand{\Ind}{{\operatorname{Ind}}}
\newcommand{\Inn}{{\operatorname{Inn}}}
\newcommand{\Int}{{\operatorname{int}}}
\newcommand{\Irr}{{\operatorname{Irr}}}
\newcommand{\Jac}{{\operatorname{Jac}}}
\newcommand{\Mor}[1]{\operatorname{Mor}(\textbf{#1})}
\newcommand{\Nil}{{\operatorname{Nil}}}
\newcommand{\Obj}{{\operatorname{Obj}}}
\newcommand{\Out}{{\operatorname{Out}}}
\newcommand{\Pic}{{\operatorname{Pic}\,}}
\newcommand{\Rep}{{\operatorname{Rep}}}
\newcommand{\Res}{{\operatorname{Res}}}
\newcommand{\Spl}{{\operatorname{Spl}}}
\newcommand{\Syl}{{\operatorname{Syl}}}
\newcommand{\Sym}{{\operatorname{Sym}}}
\newcommand{\Tor}{{\operatorname{Tor}}}

\newcommand{\Char}{{\operatorname{char}}}
\newcommand{\Frac}{{\operatorname{Frac}}}
\newcommand{\Frob}{{\operatorname{Frob}}}
\newcommand{\Proj}{{\operatorname{Proj}\,}}
\newcommand{\RMod}{{}_R\operatorname{Mod}}
\newcommand{\SExt}{{\mathcal{E}\emph{xt}}}
\newcommand{\SHom}{{\emph{Hom}}}
\newcommand{\Span}{{\operatorname{span}}}
\newcommand{\Spec}{{\operatorname{Spec}\,}}
\newcommand{\Supp}{{\operatorname{Supp}}}
\newcommand{\Tens}[1]{#1\otimes --}
\newcommand{\fHom}[1]{\operatorname{Hom}(#1, --)}

\newcommand{\mSpec}{{\operatorname{mSpec}}}


% connectors
% \renewcommand{\mid}{\big|}
\newcommand{\sse}{\subseteq}
\newcommand{\subg}{\leqslant}
\newcommand{\supg}{\geqslant}
\newcommand{\ssne}{\subsetneq}
\newcommand{\isosubg}{\lesssim}
\newcommand{\nsse}{\not\subseteq}
\newcommand{\nsubg}{\trianglelefteq}
\newcommand{\nsupg}{\trianglerighteq}
\newcommand{\csubg}{\blacktriangleleft}
\renewcommand{\mid}{\big|}


% Arrows
\newcommand{\rw}{\rightarrow}
\newcommand{\Rw}{\Rightarrow}
\newcommand{\lw}{\leftarrow}
\newcommand{\Lw}{\Leftarrow}
\newcommand{\lrw}{\leftrightarrow}
\newcommand{\LRw}{\Leftrightarrow}
\newcommand{\acton}{\curvearrowright}
\newcommand{\actson}{\curvearrowright}
\newcommand\mapsfrom{\mathrel{\reflectbox{\ensuremath{\mapsto}}}}

% arrows for commutative diagram: create four new angled double-struck arrows
\newcommand\myrot[1]{\mathrel{\rotatebox[origin=c]{#1}{$\Rightarrow$}}}
\newcommand\NEarrow{\myrot{45}}
\newcommand\NWarrow{\myrot{135}}
\newcommand\SWarrow{\myrot{-135}}
\newcommand\SEarrow{\myrot{-45}}


%shorthands
\newcommand{\oln}[1]{{\overline{#1}}}
\renewcommand{\bf}[1]{\textbf{#1}}
\newcommand{\qn}{\quad\newline}


% limits
\newcommand{\Lim}{{\varprojlim}}
\newcommand{\coLim}{{\varinjlim}}
\newcommand{\Dlim}{\lim\limits_{\longrightarrow}}
\newcommand{\coDlim}{\lim\limits_{\longleftarrow}}
\newcommand{\cofHom}[1]{\operatorname{Hom}(--, #1)}


%for saturated ideals in the localization section (I don't like these, I think I'll delete)
\newcommand{\LIm}{{}^e }
\newcommand{\LPre}{{}^c }

%  ╔══════════════════════════════════════════════════════════╗
%  ║                    Custom Commands                       ║
%  ╚══════════════════════════════════════════════════════════╝

% swap varphi and phi
\let\temp\phi
\let\phi\varphi
\let\varphi\temp

\newcommand{\placeholder}{\_\_\_\_\_}

\newcommand{\set}[2]{\left\{#1 \ \middle|\ #2\right\}}

%mod should not have the crazy amount of space to its right!
\renewcommand{\mod}{\,\operatorname{mod}\,}

% dcups
\newcommand{\dcup}{\sqcup}
\newcommand{\Dcup}{\coprod}
% \newcommand{\Dcup}{\bigsqcup}

\newcommand{\nbot}{\mathrel{\rlap{$\bot$}\mkern2mu/}}

%a nice large circle
\newcommand{\tikzcircle}[2][red,fill=black]{\tikz[baseline=-0.5ex]\draw[#1,radius=#2] (0,0) circle ;}%

\makeatletter
\newcommand*\bigcdot{\mathpalette\bigcdot@{.5}}
\newcommand*\bigcdot@[2]{\mathbin{\vcenter{\hbox{\scalebox{#2}{$\m@th#1\bullet$}}}}}
\makeatother

%somehow not a default command
\def\rddots{\mathstrut^{.^{.^{.}}}}

%% new delimiter
\DeclarePairedDelimiter{\ceil}{\lceil}{\rceil}


%  ╔══════════════════════════════════════════════════════════╗
%  ║                     environments                         ║
%  ╚══════════════════════════════════════════════════════════╝


%remember the option: breakable

% create tcolor boxes
\newtcbtheorem[number within=section]{thm}{Theorem}%
{
colback=green!5,
colframe=green!35!black,
fonttitle=\bfseries,
label type=theorem
}{th}

\newtcbtheorem[number within=section, use counter from=thm]{lem}{Lemma}%
{
colback=blue!5,
colframe=black!35!black,
fonttitle=\bfseries,
label type=lemma
}{lm}
\newtcbtheorem[number within=section, use counter from=thm]{prop}{Proposition}%
{
colback=white!5,
colframe=white!35!black,
fonttitle=\bfseries,
label type=proposition
}{pr}
\newtcbtheorem[number within=section, use counter from=thm]{cor}{Corollary}%
{colback=blue!5,
colframe=blue!35!black,
fonttitle=\bfseries,
label type=corollary
}{co}
\newtcbtheorem[number within=section, use counter from=thm]{axiom}{Axiom}%
{colback=black!5,
colframe=yellow!35!black,
fonttitle=\bfseries,
label type=axiom
}{ax}
\newtcbtheorem[number within=section, use counter from=thm]{defn}{Definition}%
{colback=red!5,
colframe=red!35!black,
fonttitle=\bfseries,
label type=definition
}{df}


% for <s-cr> to create \item[$\LRw$] instead of \item
\newenvironment{equivEnumerate}
 {\begin{enumerate}
	}{
\end{enumerate}}

%insure the exercises are properly numbered
\counterwithin{Exercise}{section}
\counterwithin{Answer}{section}


%Custom enviroments
 \newenvironment{note}
 {\begin{tcolorbox}[enhanced, sharp corners, colback=white , borderline={0.2pt}{0pt}{black}]
   \begin{quote}\textbf{Note}:
   }{
   \end{quote}\end{tcolorbox}}

 \newenvironment{intuition}
 {\begin{quote}\textbf{Intuition}:
   }{
   \end{quote}}

\newtcbtheorem[number within=section, use counter from=thm]{titledBox}{}{%
  enhanced,
  sharp corners,
  colback=white,
  colframe=black,
  borderline={0.2pt}{0pt}{black},
  left=2mm,
  right=2mm,
  top=2mm,
  bottom=2mm,
  title style={opacity=1},
  attach title to upper,
  before upper={\textbf{\tcbtitletext}\quad},
  coltitle=black,
  fonttitle=\bfseries,
  separator sign={\quad}
}{box}


%better example and proof environments
\def\exampletext{Example} % If English
\colorlet{colexam}{red!55!black} % Global example color


\newtcbtheorem[number within=section, use counter from=thm]{example}{Example}{% \newtcolorbox[use counter=testexample]{testexamplebox}{%
% Example Frame Start
empty,% Empty previously set parameters
title={\exampletext \thetcbcounter #1},% use \thetcbcounter to access the testexample counter text
% Attaching a box requires an overlay
attach boxed title to top left,
   % Ensures proper line breaking in longer titles
   minipage boxed title,
% (boxed title style requires an overlay)
boxed title style={empty,size=minimal,toprule=0pt,top=4pt,left=3mm,overlay={}},
coltitle=colexam,fonttitle=\bfseries,
before=\par\medskip\noindent,parbox=false,boxsep=0pt,left=3mm,right=0mm,top=2pt,breakable,pad at break=0mm,
   before upper=\csname @totalleftmargin\endcsname0pt, % Use instead of parbox=true. This ensures parskip is inherited by box.
% Handles box when it exists on one page only
overlay unbroken={\draw[colexam,line width=.5pt] ([xshift=-0pt]title.north west) -- ([xshift=-0pt]frame.south west); },
% Handles multipage box: first page
overlay first={\draw[colexam,line width=.5pt] ([xshift=-0pt]title.north west) -- ([xshift=-0pt]frame.south west); },
% Handles multipage box: middle page
overlay middle={\draw[colexam,line width=.5pt] ([xshift=-0pt]frame.north west) -- ([xshift=-0pt]frame.south west); },
% Handles multipage box: last page
overlay last={\draw[colexam,line width=.5pt] ([xshift=-0pt]frame.north west) -- ([xshift=-0pt]frame.south west); }%
}{ex}



\newtcolorbox{ProofTcb}[1][]{%
    empty,
    title={\emph{Proof}: #1},
    attach boxed title to top left,
    minipage boxed title,
    boxed title style={empty,size=minimal,toprule=0pt,top=4pt,left=3mm,overlay={}},
    coltitle=black!55!black,
    fonttitle=\bfseries,
    before=\par\medskip\noindent,
    parbox=false,
    boxsep=0pt,
    left=3mm,
    right=0mm,
    top=2pt,
    breakable,
    pad at break=0mm,
    before upper=\csname @totalleftmargin\endcsname0pt,
    width=\dimexpr\textwidth-3mm\relax,
    overlay unbroken={\draw[black!55!black,line width=.5pt] ([xshift=-0pt]title.north west) -- ([xshift=-0pt]frame.south west); },
    overlay first={\draw[black!55!black,line width=.5pt] ([xshift=-0pt]title.north west) -- ([xshift=-0pt]frame.south west); },
    overlay middle={\draw[black!55!black,line width=.5pt] ([xshift=-0pt]frame.north west) -- ([xshift=-0pt]frame.south west); },
    overlay last={\draw[black!55!black,line width=.5pt] ([xshift=-0pt]frame.north west) -- ([xshift=-0pt]frame.south west); },
    #1
}

\newenvironment{Proof}[1][]
  {\begin{ProofTcb}[#1]}
  {\end{ProofTcb}}


%  ╔══════════════════════════════════════════════════════════╗
%  ║                    for Dynkin diagrams                   ║
%  ╚══════════════════════════════════════════════════════════╝

\def\row#1/#2!{#1_{\IfStrEq{#2}{}{n}{#2}} & \dynkin{#1}{#2}\\}
\newcommand{\tble}[1]{
   \renewcommand*\do[1]{\row##1!}
   \[
      \begin{array}{ll}\docsvlist{#1}\end{array}
   \]
}

%  ╔══════════════════════════════════════════════════════════╗
%  ║                         links                            ║
%  ╚══════════════════════════════════════════════════════════╝

\definecolor{LinkColour}{HTML}{A64A3B} % favourite
% \definecolor{LinkColour}{HTML}{8C3D32} % 2nd place
\hypersetup{
  colorlinks,
  citecolor=black,
  filecolor=magenta,
  linkcolor=LinkColour,
  urlcolor=blue,
}



%  ╔══════════════════════════════════════════════════════════╗
%  ║                          misc                            ║
%  ╚══════════════════════════════════════════════════════════╝

\stackMath %to be able to stack text in math mode
\pgfplotsset{compat=1.18}
\makeindex

%import options: all, bibliography, book-formatting, booksSetting, customEnvironments, dynkin, hw-formatting, language-setup, newcommands, packages, preamble, proofs, settings, tcolorbox, test.svg,
\begin{document}
\maketitle
\tableofcontents

A group $G$ is \emph{[group] automorphism-representable} if there exists a group $H$ such that
\[
  \Aut_\Grp(H) \cong G
\]
From now on we'll just write $\Aut$ instead of $\Aut_\Grp$.

\begin{lem}{Rational Are Not Automorphism-representable}{rationalAreNotAutomorphism-representable}
  The rational numbers under addition $(\Q, +)$ are not automorphism representable
\end{lem}

\begin{Proof}
  The key to remember is that finitely generated subgroups of $\Q$ are cyclic.
  For the sake of contradiction say they are so that there exists a $G$ such that:
  \[
    \Aut(G) \cong \Q
  \]
  Then as:
  \[
    \Inn(G) \nsubg \Aut(G) \cong \Q
  \]
  either $\Inn(G)$ is a divisible group or cyclic. If it is is cyclic then it must have infinite order as $\Q$
  is torsion free. Then as:
  \[
    G/Z(G) \cong \Z
  \]
  Elementary group theory tells us $G$ is abelian. But if $G$ is abelian there must be a torsion element
  $x \mapsto x^{-1}$ of order $2$, a contradiction. On the other hand, if $G/Z(G)$ is not finitely generated,
  we can ``cheat'' by choosing any $gZ(G), hZ(G) \in G/Z(G)$ and generating the group $H$ which is cyclic, so
  \[
    \langle xZ(G) \rangle = \langle gZ(G), hZ(G) \rangle
  \]
  Then it is easy to see that:
  \[
    gh = x^nz_1x^mz_2 = z^{n+m}z_1z_2 = z^mz_2x^nz_1 = hg
  \]
  and as $g,h$ were arbitrary, $G$ is abelian giving us the same contradiction.
\end{Proof}

Another important fact: let us make sure that if $G$ is an $\F_2$-vector-space then $\Aut(G)$ cannot be
abelian. The point of this is that we want to abuse the $x \mapsto x^{-1}$ argument, but we can't do that if
$x^{-1} = x$ and hence this is the identity:

\begin{lem}{Vector Space, then Non-abelian}{vecSpthenNon-abelian}
  Assume $x \mapsto x^{-1}$ is the identity. Then $G$ has an $\F_2$-action, hence is a vector space and hence
  $\Aut(G)$ cannot be abelian
\end{lem}

\begin{Proof}
  Say that $x \mapsto x^{-1}$ is the identity. Then every element has a natural $\F_2$-action on it. Hence $G$
  is a vector-space which we know by linear algebra $\Aut(G)$ is not abelian.
\end{Proof}

Hence, it truly is enough to check for $x \mapsto x^{-1}$. Let's label this $\sigma$.

\subsection*{Central automorphism and cohomology}

\begin{defn}{Central Automorphisms}{centralAutos}
    An automorphism $\alpha \in \Aut(G)$ is called \emph{central} if it induces the identity map on the
    central quotient $G/Z(G)$. Equivalently, for all $x \in G$, $x^{-1}\alpha(x) \in Z(G)$. In other words:
    \[
      \alpha(x) \in xZ(G)
    \]
    so $\alpha$ keeps elements in the cosets.
\end{defn}

If we are trying to show $\Aut(G) \not\cong \Q^n$ or $\R$, note that every automorphism \emph{has} to be
central, that is:
\[
  \Aut_c(G) \cong \Aut(G)
\]
To show this, since $\Aut(G)$ is abelian, the subgroup of inner automorphisms, $\Inn(G) \cong G/Z(G)$, must be
a normal subgroup of $\Aut(G)$. Let $\alpha \in \Aut(G)$ and $c_g \in \Inn(G)$ (where $c_g(x) = gxg^{-1}$).
Since $\Aut(G)$ is abelian, $\alpha$ and $c_g$ commute:
\[
    \alpha \circ c_g = c_g \circ \alpha
\]
Applying this to all $x \in G$:
\begin{align*}
    \alpha(gxg^{-1}) &= g\alpha(x)g^{-1} \\
    \alpha(g)\alpha(x)\alpha(g)^{-1} &= g\alpha(x)g^{-1}
\end{align*}
Since $\alpha$ is surjective, as $x$ varies over $G$, $\alpha(x)$ varies over all of $G$.  Letting $y
= \alpha(x)$, re-arranging we see that:
\[
  g^{-1} \alpha(g) \, y = y \, g^{-1} \alpha(g)
\]
Letting $z = g^{-1}\alpha(g)$ we get:
\[
  zy = yz
\]
Hence, as $y$ varies through all elements of $G$,
\[
  z = g^{-1}\alpha(g) \in Z(G) \qquad \text{i.e.} \qquad \alpha(g) \in gZ(G)
\]
which is exactly the condition of a central automorphism. Thus if $\Aut(G)$ is abelian, then
\[
  \Aut(G) = \Aut_c(G)
\]
Now, each $\alpha \in \Aut_c(G)$ gives rise to a derivation. In particular, the map $\tau_\alpha: G \to Z(G)$ defined by
\[
  \tau_\alpha(x) = x^{-1}\alpha(x)
\]
is a derivation (i.e. 1-cocycle) which is easy to check as $Z(G)$ is abelian and hence these are even
homomorphisms (see~\cite{NateGaloisCohom} for more details). Since $Z(G)$ is abelian and the action is
trivial, this derivation is simply a homomorphism from $G$ to $Z(G)$. Since the target is abelian,
it factors through the commutator subgroup $G'$. Thus, we have an embedding:
\[
    \Aut_c(G) \hookrightarrow \Hom(G/G', Z(G))
\]
The condition $\Aut(G) = \Aut_c(G)$ implies that \emph{every} automorphism of $G$ acts trivially on the
quotient $Q = G/Z(G)$. This gives an interesting short exact sequence:
\[
  0 \to \Hom(G/G', Z(G)) \to \Aut_c(G) \to \Inn(G) \cap Z(\Aut(G))
\]



Now, let $Q = G/Z(G)$. This notation is used as when we were working with $\Aut(G) = \Q$ we had that $Q$ was
locally cyclic. In our case, this need not be the case, but it is still divisible and hence $Q$ is suggestive
of the useful qualities it posses. As we are assuming $G$ is non-abelian, $G/Z(G)$ is non-trivial and:
\[
  \omega: Q\times Q \to Z(G) \qquad (a,b) \mapsto [a,b]
\]
is non-trivial (if it were trivial, then for any choice of representative $xZ(G),yZ(G)$ of cosets we would get that
$[x,y] = e$ which would give that $G$ is abelian).

\chapter{Torsion-Free Automorphism Groups}

In this chapter, we explore a fascinating rigidity theorem concerning the structure of groups with "large" but abelian automorphism groups. Specifically, we investigate whether the automorphism group of a group $G$ can be isomorphic to a torsion-free abelian group like the additive group of rational numbers $\Q$, the real numbers $\R$, or products thereof. The result is a striking non-existence theorem: no such group $G$ exists. This proof weaves together basic group theory with concepts from homological algebra and the theory of nilpotent groups, highlighting the tension between the commutativity of $\Aut(G)$ and the non-commutative structure required of $G$.

\section{The Main Theorem}

The central question is simple to state but requires careful analysis to answer: Does there exist a group $G$ such that $\Aut(G) \cong A$, where $A$ is a torsion-free abelian group? The answer is negative.

\begin{thm}{Non-Existence of Groups with Torsion-Free Abelian Automorphisms}{nonExistAuto}
    Let $A$ be a non-trivial torsion-free abelian group (e.g., $A = \Q$, $A = \R$, $A = \Q \times \Q$). There does not exist a group $G$ such that $\Aut(G) \cong A$.
\end{thm}

The proof naturally divides into two cases based on the commutativity of $G$: the abelian case and the non-abelian case. Each case leads to a contradiction arising from the specific structural properties of $A$.

\section{The Abelian Case}

Let us first assume that $G$ is abelian. In this setting, we can leverage the existence of a specific, well-behaved automorphism: the inversion map.

\subsection{The Inversion Trap}

Consider the map $\phi: G \to G$ defined by $\phi(x) = x^{-1}$. Since $G$ is abelian, this map is a homomorphism:
\[
    \phi(xy) = (xy)^{-1} = y^{-1}x^{-1} = x^{-1}y^{-1} = \phi(x)\phi(y)
\]
Since every element in a group has a unique inverse, $\phi$ is a bijection, and thus $\phi \in \Aut(G)$.
Observe that $\phi^2(x) = (x^{-1})^{-1} = x$, so $\phi^2 = \id_G$. This means $\phi$ is an element of order
$2$ in $\Aut(G)$ (unless $\phi = \id_G$, in which case its order is $1$).

However, we assumed $\Aut(G) \cong A$, and $A$ is \emph{torsion-free}. A torsion-free group contains no
non-identity elements of finite order. Therefore, we must have $\phi = \id_G$. This implies:
\[
    x^{-1} = x \quad\iff\quad x^2 = e \quad \forall x \in G
\]
A group satisfying $x^2 = e$ for all elements is known as an \emph{elementary abelian 2-group}. Such a group
carries the structure of a vector space over the field $\F_2$.

\subsection{The Vector Space Barrier}

If $G$ is an elementary abelian 2-group, it is isomorphic to a direct sum of copies of $\Z_2$. Let $V$ denote this vector space. The automorphisms of $G$ are precisely the invertible linear transformations of $V$. Thus:
\[
    \Aut(G) \cong \GL(V)
\]
We now analyze the structure of the General Linear Group $\GL(V)$:
\begin{enumerate}
    \item If $\dim(V) = 0$ or $1$ (i.e., $|G| = 1$ or $2$), then $\Aut(G)$ is the trivial group. This contradicts the assumption that $A$ is non-trivial (since $A$ is isomorphic to $\Q$, $\R$, etc.).
    \item If $\dim(V) \ge 2$, then $\GL(V)$ is \emph{non-abelian}. It is a standard result of linear algebra that for dimension $\ge 2$, there exist non-commuting matrices (linear transformations).
\end{enumerate}
Since we assumed $\Aut(G) \cong A$ is abelian, we have reached a contradiction. Thus, $G$ cannot be abelian.

\section{The Non-Abelian Case}

We now turn to the more subtle case where $G$ is non-abelian. Here, the contradiction arises from the interplay between central automorphisms and the existence of "scaling" automorphisms in torsion-free nilpotent groups.

\subsection{Central Automorphisms and Cohomology}

Assume $G$ is non-abelian. Since $\Aut(G)$ is abelian, the subgroup of inner automorphisms, $\Inn(G) \cong G/Z(G)$, must be a normal  subgroup of $\Aut(G)$.
Let $\alpha \in \Aut(G)$ and $c_g \in \Inn(G)$ (where $c_g(x) = gxg^{-1}$). Since $\Aut(G)$ is abelian, $\alpha$ and $c_g$ commute:
\[
    \alpha \circ c_g = c_g \circ \alpha
\]
Applying this to an arbitrary $x \in G$:
\begin{align*}
    \alpha(gxg^{-1}) &= g\alpha(x)g^{-1} \\
    \alpha(g)\alpha(x)\alpha(g)^{-1} &= g\alpha(x)g^{-1}
\end{align*}
Since $\alpha$ is surjective, as $x$ varies over $G$, $\alpha(x)$ varies over all of $G$. Thus, for the conjugation actions to be identical, the conjugating elements must differ by an element of the center:
\[
    \alpha(g)g^{-1} \in Z(G) \quad \forall g \in G
\]
This condition defines a specific class of automorphisms known as \emph{central automorphisms}, denoted $\Aut_c(G)$.
\begin{defn}{Central Automorphisms}{centralAutos2}
    An automorphism $\alpha \in \Aut(G)$ is called \emph{central} if it induces the identity map on the central quotient $G/Z(G)$. Equivalently, for all $x \in G$, $x^{-1}\alpha(x) \in Z(G)$.
\end{defn}
We have thus proven that if $\Aut(G)$ is abelian, then $\Aut(G) = \Aut_c(G)$.


\subsection{The Scaling Contradiction}

We now demonstrate that for the groups in question, there must exist automorphisms that are \emph{not} central, contradicting the previous section.

Since $\Inn(G) \le \Aut(G)$ and $\Aut(G)$ is torsion-free abelian, the quotient $Q = G/Z(G)$ is a torsion-free abelian group. Because $G$ is non-abelian, the rank of $Q$ must be at least 2 (if it were rank 1, i.e., locally cyclic, $G$ would be abelian).
We also have a non-trivial commutator map $\omega: Q \times Q \to Z(G)$ defined by $\omega(\bar{x}, \bar{y}) = [x, y]$.

In the context of torsion-free nilpotent groups (specifically class 2, which $G$ must be), the structure is flexible enough to admit "scaling" automorphisms. For groups resembling the Heisenberg group over a field or ring that allows division (like subrings of $\R$), we can define automorphisms that act non-trivially on $Q$.

Let $\lambda \in \Q$ (or a suitable scalar depending on the specific structure of $Q$) such that $\lambda \neq 0, 1$. We can construct a map that "scales" the generators of $G$ modulo the center:
\begin{itemize}
    \item On the quotient $Q$, the map induces $\bar{x} \mapsto \lambda \bar{x}$ (using the additive notation for the abelian group $Q$).
    \item To preserve the commutator relation $[x, y]$, the center $Z(G)$ must be scaled by $\lambda^2$.
\end{itemize}
Such a map, call it $\psi_\lambda$, is an automorphism of $G$. Crucially, consider its action on the quotient $Q$:
\[
    \psi_\lambda(\bar{x}) = \lambda \bar{x} \neq \bar{x} \quad (\text{since } \lambda \neq 1)
\]
Because $\psi_\lambda$ does not induce the identity map on $G/Z(G)$, it is **not** a central automorphism.

\begin{itemize}
    \item Assumption that $\Aut(G)$ is abelian $\implies$ All automorphisms are central ($\Aut(G) = \Aut_c(G)$).
    \item Structure of torsion-free non-abelian $G$ $\implies$ Existence of non-central (scaling) automorphisms.
\end{itemize}

This is a contradiction. Thus, $G$ cannot be non-abelian.

\section{Conclusion}

We have exhausted all possibilities. If $G$ is abelian, the torsion-free property of $A$ forces $G$ to be an elementary abelian 2-group, making $\Aut(G)$ non-abelian (or trivial). If $G$ is non-abelian, the abelian property of $A$ forces rigidity ($\Aut(G) = \Aut_c(G)$) that is incompatible with the flexible scaling automorphisms allowed by the torsion-free structure.

Therefore, no group $G$ exists with $\Aut(G) \cong \Q$, $\Aut(G) \cong \R$, or any other non-trivial torsion-free abelian group.

\begin{Exercise}[label=ex:AutoGroups]
  \Question Show that if $G$ is a group of exponent 2, then $G$ is abelian.
  \Question Prove that for the Heisenberg group $H(\Z)$, $\Aut(H(\Z))$ is not abelian by exhibiting a non-central automorphism.
  \Question Verify that the inversion map $x \mapsto x^{-1}$ is an automorphism if and only if the group is abelian.
\end{Exercise}
\begin{Answer}[ref={ex:AutoGroups}]
  \Question Consider $(xy)^2 = e$. Then $xyxy = e$. Multiply by $x$ on left and $y$ on right: $yxy = x^{-1}y^{-1} = xy$. Then $yx = xy$.
\end{Answer}


\section*{Clarifications on the Proof}

\subsection*{1. The Ordering: $\alpha(g)g^{-1}$ vs $g^{-1}\alpha(g)$}

You are correct to spot this. In the derivation of the contradiction, we derived that $\alpha(g)$ and $g$ must commute modulo the center.
Technically, from $\alpha(g)y\alpha(g)^{-1} = gyg^{-1}$, we deduce that the element $z = g^{-1}\alpha(g)$ commutes with every $y \in G$, placing $z$ in the center $Z(G)$.

Why are $\alpha(g)g^{-1}$ and $g^{-1}\alpha(g)$ interchangeable here?
\begin{enumerate}
    \item Note that $\alpha(g)g^{-1} = g(g^{-1}\alpha(g))g^{-1}$.
    \item This shows that $\alpha(g)g^{-1}$ is the \emph{conjugate} of $g^{-1}\alpha(g)$ by $g$.
    \item Since we established that $g^{-1}\alpha(g)$ is in the Center $Z(G)$, and central elements are invariant under conjugation, it follows that:
    \[ \alpha(g)g^{-1} = g^{-1}\alpha(g) \]
\end{enumerate}
So, while $g^{-1}\alpha(g)$ is the standard definition for the cohomology derivation, they represent the same element in this specific context.

\subsection*{2. Why is $\tau_\alpha$ a 1-cocycle?}

We claimed that the map $\tau_\alpha(x) = x^{-1}\alpha(x)$ is a derivation (or 1-cocycle) from $G$ to $Z(G)$.
A map $f: G \to M$ (where $G$ acts on the module $M$) is a 1-cocycle if it satisfies the condition:
\[ f(xy) = f(x) + x \cdot f(y) \]
(using additive notation for the module $M$).

In our case, $M = Z(G)$. The action of $G$ on its own center is **trivial** (i.e., $g \cdot z = z$ for all $g \in G, z \in Z(G)$). Thus, the condition simplifies to the standard homomorphism property (written multiplicatively):
\[ \tau(xy) = \tau(x)\tau(y) \]
Let's verify this for $\tau_\alpha(x) = x^{-1}\alpha(x)$:
\begin{align*}
    \tau_\alpha(xy) &= (xy)^{-1}\alpha(xy) \\
    &= y^{-1}x^{-1}\alpha(x)\alpha(y)
\end{align*}
We know that $x^{-1}\alpha(x)$ is in the center $Z(G)$. Central elements commute with everything, so we can swap $x^{-1}\alpha(x)$ with $y^{-1}$:
\begin{align*}
    &= y^{-1} [x^{-1}\alpha(x)] \alpha(y) \\
    &= [x^{-1}\alpha(x)] [y^{-1}\alpha(y)] \\
    &= \tau_\alpha(x) \tau_\alpha(y)
\end{align*}
Thus, $\tau_\alpha$ is a homomorphism from $G$ to $Z(G)$, which makes it a 1-cocycle for the trivial action.

\subsection*{3. Rank of $Q$}

When we speak of the **rank** of the torsion-free abelian group $Q = G/Z(G)$, we refer to the **Torsion-Free Rank** (also known as the $\Q$-rank).
Ideally, if we view $Q$ as an additive group, the rank is the dimension of the vector space $Q \otimes_{\Z} \Q$.
\[ \text{rank}(Q) = \dim_{\Q} (Q \otimes \Q) \]
Intuitively, this is the size of the largest subset of elements in $Q$ that are linearly independent over the integers (i.e., no linear combination with integer coefficients equals 0).

\textbf{Why Rank $\ge 2$?}
If the rank of $Q$ is 1, then $Q$ is isomorphic to a subgroup of $\Q$.
As shown in the simpler case, subgroups of $\Q$ are **locally cyclic** (every finitely generated subgroup is cyclic).
If $G/Z(G)$ is locally cyclic, then $G$ must be abelian.
\begin{itemize}
    \item \emph{Sketch:} For any $x, y \in G$, their projections $\bar{x}, \bar{y}$ in a locally cyclic quotient generate a cyclic subgroup. Thus $\bar{x} = \bar{g}^n, \bar{y} = \bar{g}^m$. This implies $x$ and $y$ are powers of the same element $g$ modulo the center, so they commute.
\end{itemize}
Since we are in the **Non-Abelian Case**, $Q$ cannot be rank 1. Thus, $\text{rank}(Q) \ge 2$.

\subsection*{4. Why is the Commutator Non-Trivial?}

The commutator map $\omega: Q \times Q \to Z(G)$ is defined by $\omega(\bar{x}, \bar{y}) = xyx^{-1}y^{-1}$.
If this map were trivial (i.e., the image is always the identity $e$), then:
\[ xyx^{-1}y^{-1} = e \implies xy = yx \quad \forall x,y \in G \]
This would imply $G$ is abelian.
Since we assume $G$ is **non-abelian**, there must exist at least one pair $x, y$ that does not commute, making the commutator map non-trivial.

\subsection*{5. The Existence of $\psi_\lambda$}

This is the most subtle point. We cannot define $x^\lambda$ (for $\lambda \in \Q$) in just \emph{any} group. However, the assumption that $\Aut(G)$ is a large torsion-free divisible group (like $\Q$ or $\R$) implies strong structural properties on $G$.

Specifically, $G$ must be a subgroup of a **divisible nilpotent group** (specifically, its \emph{Mal'cev Completion}). Divisible nilpotent groups behave exactly like Lie algebras over $\Q$. In such groups, we can define rational powers $g^\lambda$.

Let's look at a concrete example using the simplest non-abelian group of this type: the **Heisenberg Group** $H$.
\[ H = \langle x, y, z \mid [x, y] = z, [x, z] = [y, z] = 1 \rangle \]
Here, $Z(G) = \langle z \rangle$ and $Q = G/Z(G) \cong \Z \times \Z$ (generated by $\bar{x}, \bar{y}$).
We want to define an automorphism that scales $x$ but not $y$.
Let $\lambda = 2$.
We cannot simply say $\psi(x) = x^2$ and $\psi(y) = y$ because that might break the commutator relation.
Let's verify the commutator relation $[x^2, y]$.
Using the identity $[a^n, b] = [a, b]^n$ (which holds when $[a,b]$ is central):
\[ [x^2, y] = [x, y]^2 = z^2 \]
So, if we define the map as:
\[ \psi(x) = x^2, \quad \psi(y) = y, \quad \psi(z) = z^2 \]
This preserves the defining relation $\psi([x, y]) = \psi(z) = z^2$ and $[\psi(x), \psi(y)] = [x^2, y] = z^2$.
Thus, $\psi$ is a valid endomorphism. In the context of divisible groups (where we can take roots), we can inverse this with $\lambda = 1/2$, so it is an automorphism.

\textbf{The Contradiction:}
Is $\psi$ a central automorphism?
Check its action on the quotient $Q$:
\[ \psi(\bar{x}) = \overline{x^2} = 2\bar{x} \neq \bar{x} \]
Since $\psi$ acts non-trivially on the quotient (it doubles the generator $\bar{x}$), it is **not** a central automorphism.
This contradicts the earlier deduction that $\Aut(G) \text{ abelian} \implies \Aut(G) = \Aut_c(G)$.


\chapter{lie algebra stuff}

\section{Step 2: The Lie Algebra Connection}

To find a contradiction, we must show that there exists an automorphism that is \emph{not} central. To do this, we need to understand the structure of $G$.
Since $\Aut(G) \cong \Q$ (or $\R$), $G$ must be a divisible group (we can take $n$-th roots of elements).
This invokes a powerful correspondence in algebra.

\begin{thm}{Mal'cev Correspondence}{malcevCorr}
    There is a bijective correspondence between divisible torsion-free nilpotent groups and rational nilpotent Lie algebras.
\end{thm}

For our purposes, this means we can study $G$ by studying a corresponding Lie Algebra $\mathfrak{g}$ over the field $\Q$.
\begin{itemize}
    \item The Lie Algebra $\mathfrak{g}$ is a vector space over $\Q$ equipped with a bracket $[X, Y]$.
    \item There is a bijective map $\exp: \mathfrak{g} \to G$. Unlike in rotation groups (where $\exp$ wraps around), for nilpotent groups this map is a **global bijection** (diffeomorphism).
    \item Every element $g \in G$ has a unique logarithm $X \in \mathfrak{g}$ such that $g = \exp(X)$.
\end{itemize}

\subsubsection{The Campbell-Baker-Hausdorff (CBH) Formula}
You asked how the multiplication in $G$ is determined by the algebra. The answer is the **CBH Formula**:
\[ \exp(X) \cdot \exp(Y) = \exp\left( X + Y + \frac{1}{2}[X, Y] + \frac{1}{12}[X, [X, Y]] - \frac{1}{12}[Y, [X, Y]] + \dots \right) \]
Since $\Aut(G)$ is abelian, $G$ is nilpotent of **class 2**. This means all commutators of commutators vanish ($[X, [Y, Z]] = 0$).
Consequently, the infinite series truncates at the second term:
\[ \exp(X) \cdot \exp(Y) = \exp\left( X + Y + \frac{1}{2}[X, Y] \right) \]
This formula tells us that group multiplication is just vector addition plus a "twist" given by the commutator.

\subsection{Step 3: Constructing the Scaling Automorphism}

We can now define coordinates. Since $\mathfrak{g}$ is a vector space, let $\{X, Y\}$ be basis vectors for the quotient $\mathfrak{g}/Z(\mathfrak{g})$ and $Z = [X, Y]$ be the basis for the center.
Any $g \in G$ corresponds to a unique vector $v = aX + bY + cZ$. We can represent $g$ as $(a, b, c)$.

We define a scaling map $\psi_\lambda$ for some scalar $\lambda \in \Q$ ($\lambda \neq 0, 1$).
In the Lie Algebra (vector space), scaling is easy. We scale the generators $X$ and $Y$ by $\lambda$.
\[ X \mapsto \lambda X, \quad Y \mapsto \lambda Y \]
What happens to the center vector $Z$? Since $Z = [X, Y]$ and the bracket is bilinear:
\[ [\lambda X, \lambda Y] = \lambda^2 [X, Y] = \lambda^2 Z \]
So we must scale the center component by $\lambda^2$.

Let us lift this map to the group $G$. Define $\psi_\lambda: G \to G$ by applying this scaling to the coordinates:
\[ \psi_\lambda( \exp(aX + bY + cZ) ) = \exp( \lambda a X + \lambda b Y + \lambda^2 c Z ) \]

\textbf{Is this a group homomorphism?}
We verify using the truncated CBH formula. Let $g_1 \leftrightarrow A$ and $g_2 \leftrightarrow B$.
\[ \psi_\lambda(g_1 \cdot g_2) = \psi_\lambda(\exp(A + B + \frac{1}{2}[A, B])) \]
Since $\psi_\lambda$ is linear on vector components and quadratic on the bracket term:
\[ = \exp\left( \lambda(A+B) + \lambda^2 \left(\frac{1}{2}[A, B]\right) \right) \]
(Note: The bracket $[A, B]$ lies in the center, so it scales by $\lambda^2$).

Now compute $\psi_\lambda(g_1) \cdot \psi_\lambda(g_2)$. Let $A' = \lambda A$ (on generators) and $B' = \lambda B$.
\[ \psi_\lambda(g_1) \psi_\lambda(g_2) = \exp(A') \cdot \exp(B') = \exp\left( A' + B' + \frac{1}{2}[A', B'] \right) \]
Substituting the scaling:
\[ = \exp\left( \lambda A + \lambda B + \frac{1}{2}[\lambda A, \lambda B] \right) = \exp\left( \lambda(A+B) + \frac{1}{2}\lambda^2 [A, B] \right) \]
The results match. Thus, $\psi_\lambda$ is a homomorphism. Since $\lambda \neq 0$ and $G$ is divisible, it is invertible, hence an automorphism.

\subsection{Step 4: The Contradiction}

We check if $\psi_\lambda$ is a central automorphism.
Consider the generator $g = \exp(X)$.
\[ \psi_\lambda(g) = \exp(\lambda X) = g^\lambda \]
In the quotient group $G/Z(G)$, this maps the coset $\bar{g}$ to $\bar{g}^\lambda$.
If $\lambda \neq 1$, then $\bar{g}^\lambda \neq \bar{g}$.
Therefore, $\psi_\lambda$ acts non-trivially on the quotient. It is **not** a central automorphism.

\textbf{Conclusion:}
\begin{itemize}
    \item Assumption ($\Aut(G)$ abelian) $\implies$ All automorphisms are central.
    \item Structure (Divisible Nilpotent) $\implies$ There exist non-central scaling automorphisms.
\end{itemize}
This is a contradiction. Thus, no such group $G$ exists.

\chapter{another idea}



We prove that there does not exist a group $G$ such that $\operatorname{Aut}(G)$ is isomorphic to a subgroup of $\mathbb{Q}^n$ (or any non-trivial torsion-free divisible abelian group). The proof relies on analyzing the structure of the inner automorphism group and constructing a specific "inversion" automorphism that leads to a contradiction.

\section{Introduction}

Let $A$ be a non-trivial torsion-free divisible abelian group (e.g., $A = \mathbb{Q}$, $A = \mathbb{R}$, or $A = \mathbb{Q}^n$). We aim to prove that there is no group $G$ such that $\operatorname{Aut}(G) \cong A$.

\section{Proof}

\begin{thm}{main theorem}{mainthm}
There is no group $G$ such that $\operatorname{Aut}(G) \cong A$, where $A$ is a subgroup of $\mathbb{Q}^n$ containing at least one non-zero element.
\end{thm}

\begin{proof}
Suppose such a group $G$ exists. Let $\mathcal{A} = \operatorname{Aut}(G) \cong A$.

\subsection*{Step 1: Structural Consequences for $G$}
Since $\mathcal{A}$ is abelian, every subgroup of $\mathcal{A}$ is normal. In particular, the group of inner automorphisms, $\operatorname{Inn}(G) \cong G/Z(G)$, is a normal subgroup of $\mathcal{A}$.
\begin{enumerate}
    \item Since $\mathcal{A}$ is abelian, $\operatorname{Inn}(G)$ is abelian. Thus, $G/Z(G)$ is abelian, which implies that $G$ is nilpotent of class at most 2. The commutator subgroup $G'$ is contained in the center $Z(G)$.
    \item Since $\mathcal{A}$ is torsion-free (being a subgroup of $\mathbb{Q}^n$), $\operatorname{Inn}(G)$ is torsion-free. Thus, the quotient $Q = G/Z(G)$ is a torsion-free abelian group.
\end{enumerate}

\subsection*{Step 2: Constraint on Automorphisms}
Since $\mathcal{A}$ is abelian, every automorphism $\alpha \in \mathcal{A}$ must commute with every inner automorphism $c_g \in \operatorname{Inn}(G)$ (where $c_g(x) = gxg^{-1}$).
\[
\alpha \circ c_g = c_g \circ \alpha
\]
Applying this to an arbitrary element $x \in G$:
\begin{align*}
\alpha(gxg^{-1}) &= g\alpha(x)g^{-1} \\
\alpha(g)\alpha(x)\alpha(g)^{-1} &= g\alpha(x)g^{-1}
\end{align*}
Since $\alpha$ is surjective, let $y = \alpha(x)$. As $x$ ranges over $G$, $y$ ranges over $G$.
\[
\alpha(g) y \alpha(g)^{-1} = g y g^{-1} \quad \forall y \in G
\]
Multiplying by $g^{-1}$ on the left and $\alpha(g)$ on the right:
\[
(g^{-1}\alpha(g)) y = y (g^{-1}\alpha(g))
\]
This implies that the element $g^{-1}\alpha(g)$ commutes with every element $y \in G$. Therefore, $g^{-1}\alpha(g) \in Z(G)$, or equivalently:
\[
\alpha(g) \equiv g \pmod{Z(G)} \quad \forall g \in G
\]
Thus, every automorphism of $G$ is a \textbf{central automorphism}; it induces the identity map on the quotient group $G/Z(G)$.

\subsection*{Step 3: Construction of the Inversion Automorphism $\Psi$}
We define a specific map $\Psi: G \to G$ and show it is an automorphism that violates the constraint established in Step 2.

Let $Q = G/Z(G)$ and $Z = Z(G)$. Since $G$ is a central extension of $Z$ by $Q$, elements of $G$ can be identified with pairs $(a, z)$ where $a \in Q$ and $z \in Z$, with the multiplication rule:
\[
(a, z_1) \cdot (b, z_2) = (a+b, z_1 + z_2 + f(a, b))
\]
where $f: Q \times Q \to Z$ is a 2-cocycle (factor set). (We use additive notation for the abelian groups $Q$ and $Z$).

According to the theory of nilpotent groups (see Warfield, \textit{Nilpotent Groups}, Theorem 5.4), for torsion-free groups, the factor set $f$ corresponds to an alternating bilinear form (the commutator map). Bilinearity implies:
\[
f(-a, -b) = (-1)(-1)f(a, b) = f(a, b)
\]
We define the map $\Psi: G \to G$ by:
\[
\Psi((a, z)) = (-a, z)
\]
This map inverts elements in the quotient $Q$ and fixes elements in the center $Z$.

\textbf{Verification that $\Psi$ is a Homomorphism:}
\begin{align*}
\Psi((a, z_1) \cdot (b, z_2)) &= \Psi(a+b, z_1+z_2+f(a, b)) \\
&= (-(a+b), z_1+z_2+f(a, b))
\end{align*}
\begin{align*}
\Psi(a, z_1) \cdot \Psi(b, z_2) &= (-a, z_1) \cdot (-b, z_2) \\
&= (-a-b, z_1+z_2+f(-a, -b))
\end{align*}
Since $f(a, b) = f(-a, -b)$, the map preserves the group operation. Since $\Psi^2 = \text{id}$, it is bijective and thus an automorphism.

\subsection*{Step 4: The Contradiction}
We have established two contradictory facts about $\Psi$:
\begin{enumerate}
    \item \textbf{From Step 2:} Since $\Psi \in \operatorname{Aut}(G)$ and $\operatorname{Aut}(G)$ is abelian, $\Psi$ must be a central automorphism. It must induce the identity map on $Q$.
    \[ \Psi(a, z) \equiv (a, z) \pmod Z \implies -a = a \]
    \item \textbf{From Step 3:} By definition, $\Psi$ induces the inversion map on $Q$.
    \[ \Psi(a, z) = (-a, z) \]
\end{enumerate}
These conditions require $-a = a$, or $2a = 0$, for all $a \in Q$.
Since $Q$ is a subgroup of $\mathbb{Q}^n$, it is **torsion-free**. The only solution to $2a=0$ is $a=0$.
This implies $Q$ is the trivial group, so $G = Z(G)$, and $G$ is abelian.

\textbf{The Abelian Case:}
If $G$ is abelian, the map $\phi: x \mapsto x^{-1}$ is an automorphism.
\[ \phi(x^{-1}) = (x^{-1})^{-1} = x \implies \phi^2 = \text{id} \]
Thus $\operatorname{Aut}(G)$ contains an element of order 2.
But $\operatorname{Aut}(G) \cong A \le \mathbb{Q}^n$, which is torsion-free.
Therefore $\phi = \text{id}$, which implies $x^{-1} = x$ for all $x$.
$G$ must be an elementary abelian 2-group.
The automorphism group of an elementary abelian 2-group of rank $k > 1$ is $\operatorname{GL}(k, \mathbb{F}_2)$, which is non-abelian. If rank $k=1$, $\operatorname{Aut}(G)$ is trivial.
Since $A$ is non-trivial and abelian, this is a contradiction.

\end{proof}

\chapter{another attempt}


\section{Theorem}
There does not exist a group $G$ such that $\operatorname{Aut}(G) \cong A$, where $A$ is a subgroup of $\mathbb{Q}^n$ (or any non-trivial torsion-free divisible abelian group).

\section{Proof}

Assume for the sake of contradiction that such a group $G$ exists. Let $\mathcal{A} = \operatorname{Aut}(G)$. By hypothesis, $\mathcal{A}$ is torsion-free and abelian.

\subsection*{Step 1: Structural Reduction}
Since $\mathcal{A}$ is abelian, the inner automorphism group $\operatorname{Inn}(G) \cong G/Z(G)$ is a subgroup of $\mathcal{A}$.
\begin{enumerate}
    \item $G/Z(G)$ is abelian, implying $G$ is nilpotent of class at most 2.
    \item $G/Z(G)$ is a subgroup of $\mathbb{Q}^n$, so it is a **torsion-free** abelian group.
\end{enumerate}
Let $Q = G/Z(G)$ and $Z = Z(G)$. The group $G$ is a central extension of $Z$ by $Q$.

\subsection*{Step 2: The Centrality Constraint}
Since $\mathcal{A}$ is abelian, every automorphism $\alpha \in \mathcal{A}$ must commute with every inner automorphism $c_g$.
\[ \alpha \circ c_g = c_g \circ \alpha \implies \alpha(gxg^{-1}) = g\alpha(x)g^{-1} \]
This implies $\alpha(g) y \alpha(g)^{-1} = g y g^{-1}$ for all $y$, forcing $g^{-1}\alpha(g) \in Z$.
Thus, every automorphism is \textbf{central}:
\[ \alpha(g) \equiv g \pmod{Z} \quad \forall \alpha \in \mathcal{A}, g \in G \]
In particular, every automorphism must induce the \textbf{Identity Map} on the quotient $Q$.

\subsection*{Step 3: The Factor Set and Commutator}
We construct a specific automorphism to contradict Step 2.
By Warfield (\textit{Nilpotent Groups}, Theorem 5.4), the extension class of $G$ in $H^2(Q, Z)$ maps to a commutator form $\omega \in \operatorname{Hom}(\Lambda^2 Q, Z)$.
\[ \omega(u, v) = [u, v] \]
The theorem states that $\omega$ is an **Alternating Bilinear Form**.
Bilinearity implies:
\[ \omega(-u, -v) = (-1)(-1)\omega(u, v) = \omega(u, v) \]
This signifies that the commutator relations are invariant if we invert the generators in $Q$.

\subsection*{Step 4: The Inversion Automorphism}
We define a map $\Psi: G \to G$. Representing elements of $G$ as pairs $(q, z)$ with factor set $f$:
\[ \Psi( (q, z) ) = (-q, z) \]
(Invert the quotient, fix the center).
We verify $\Psi$ is a homomorphism. The condition is $f(u, v) = f(-u, -v)$.
\begin{enumerate}
    \item The antisymmetric part of $f$ is the commutator $\omega$, which satisfies $\omega(u, v) = \omega(-u, -v)$ by bilinearity (Step 3).
    \item The symmetric part of $f$ corresponds to an abelian extension. Since $Q$ is torsion-free, we can choose a symmetric bilinear cocycle, or observe that for abelian groups, inversion is always an automorphism.
\end{enumerate}
Thus, $\Psi$ is a well-defined automorphism of $G$.

\subsection*{Step 5: The Contradiction}
We now have an automorphism $\Psi \in \operatorname{Aut}(G)$.
\begin{enumerate}
    \item \textbf{Torsion:} $\Psi^2(q, z) = (-(-q), z) = (q, z)$. So $\Psi^2 = \text{id}$.
    \item \textbf{Non-Centrality:} By definition, $\Psi$ acts as $-1$ on the quotient $Q$.
    From Step 2, all automorphisms must act as $+1$ on $Q$.
    This forces $-1 = +1$ on $Q$.
    \[ -q = q \implies 2q = 0 \]
    Since $Q$ is torsion-free, $q=0$.
    \item \textbf{Collapse:} This implies $Q$ is trivial, so $G = Z$, and $G$ is Abelian.
    \item \textbf{Abelian Case:} If $G$ is Abelian, the map $\phi(g) = g^{-1}$ is an automorphism of order 2.
    Since $\operatorname{Aut}(G) \le \mathbb{Q}^n$ is torsion-free, $\phi$ must be identity.
    $g^{-1} = g \implies g^2 = 1$.
    $G$ is an elementary abelian 2-group.
    $\operatorname{Aut}(G) \cong \operatorname{GL}(k, \mathbb{F}_2)$.
    For $k \ge 2$, this is non-abelian. For $k=1$, it is trivial.
    In neither case is it isomorphic to a non-trivial subgroup of $\mathbb{Q}^n$ (which is abelian and infinite).
\end{enumerate}
\textbf{Conclusion:} No such group exists.



\newpage
\printbibliography
\end{document}
